\section{BACKGROUND AND LITERATURE REVIEW}
\label{sec:background}

\subsection{Debate as Transformative Game}

Academic debate formats are games. Following \textcite{salen2004rules}, a game is a system in which players engage in artificial conflict, defined by rules, resulting in a quantifiable outcome. Academic debate satisfies all criteria: rules, participants, winners, losers, tournaments, and scoring systems. Alfred Snider, whose school of debate of training I was priviledge in being apprenticed in, through his 1982 doctoral dissertation at the University of Kansas, along with his subsequent publication \parencite{snider1984gaming}, established the scholarly foundation for viewing debate as a game while documenting its educational benefits.

But what makes a game \textit{transformative}? \textcite{bowman2025transformative} conceptualize transformative role-playing games as those intentionally designed to encourage personal and social change arising from play. For transformation to occur, this change must expand beyond the bounds of the original experience and integrate into the participant's daily frames of reality and identity. By this standard, a game is truly transformative not merely when participants report enjoyment or in-game learning, but when the experience produces lasting changes in cognition, behavior, or self-concept.

Forensics pedagogy documents debate's transformative potential across several domains that satisfy this criterion. The most extensively studied is critical thinking: debaters develop the ability to recognize stated and unstated assumptions, evaluate causal inferences, and assess the validity of arguments \parencite{colbert1995enhancing, hill1993value}. Beyond cognitive gains, debate serves as a vehicle for empowerment, particularly for previously ``silenced student populations'' who develop their own voice through the activity \parencite{snider2003gamemaster}. The activity cultivates what \textcite{hicks1998public} calls a ``democratic ethos'': the capacity to justify convictions in public, respect minority views, and evaluate opposing arguments critically. Oral communication proficiency develops through audience adaptation, where students learn to analyze listeners' beliefs and norms to frame compelling narratives \parencite{cirlin1997sociological}, while parliamentary formats foster extemporaneous thinking and speaking \parencite{whitmore1998british}. Debate also provides a socialization mechanism, integrating students into a sub-culture with specialized language that enhances their sense of belonging to an intellectual community \parencite{cirlin1997sociological, kargacin1998state}. Crucially, these effects transfer beyond the activity into academic performance, professional communication, and civic engagement, satisfying the criterion of transformation beyond the game.

% Empirical research substantiates these claims, with quantitative evidence accumulating across critical thinking, academic achievement, and college readiness over several decades. \textcite{allen1999meta} conducted a meta-analysis of studies examining forensics and communication education, finding a significant positive relationship between such training and gains in critical thinking (longitudinal designs $r = .18$, cross-sectional designs $r = .20$), with forensics participation demonstrating the largest effect among the communication activities studied. At the individual study level, \textcite{colbert1987effects} conducted the most comprehensive investigation, tracking 275 CEDA and NDT debaters and non-debaters across nine universities over a full academic year. Using analysis of covariance controlling for pretest scores, he found that debaters' Watson-Glaser Critical Thinking Appraisal scores rose from 61.18 to 64.53 while non-debaters' scores actually declined from 52.67 to 49.14, with NDT debaters significantly outgaining non-debaters ($F = 31.77$, $p < .001$). While \textcite{hill1993value} cautioned that the broader body of debate-critical thinking research remained inconclusive due to methodological limitations and inconsistent results across sub-samples, \textcite{colbert1995enhancing} responded that the preponderance of evidence establishes ``presumptive proof'' favoring a positive debate-critical thinking relationship, particularly when properly designed studies using robust statistical controls are considered.
%
% Beyond critical thinking, debate participation correlates with substantial academic outcomes. \textcite{mezuk2009urban} analyzed data from the Chicago Debate League spanning 1997 to 2006 and found that participants were 70\% more likely to graduate and three times less likely to drop out than comparable non-participants, even after controlling for eighth-grade test scores and grade point average. Debaters were also more likely to meet ACT benchmarks for college readiness in English and reading, though not in science or mathematics. A subsequent expanded study of the same program by \textcite{mezuk2011impact}, analyzing 9,145 students of whom 2,449 participated in at least one CDL tournament, found that debaters were 2.5 times less likely to drop out, graduated at a rate of 90\% compared to 75\% for non-debaters, and achieved a mean twelfth-grade GPA of 3.23 compared to 2.83 for their peers. Most recently, \textcite{schueler2025debate} examined policy debate in Boston Public Schools serving large concentrations of economically disadvantaged students and found positive impacts on English Language Arts scores of 0.13 standard deviations, equivalent to 68\% of a full year of average ninth-grade learning. Critically, these gains were concentrated on analytical rather than rote subskills and were largest among the lowest-performing students at baseline, with additional evidence of positive effects on high school graduation and postsecondary enrollment.

Empirical research substantiates these claims. \textcite{allen1999meta} conducted a meta-analysis finding a significant positive relationship between forensics participation and critical thinking gains (longitudinal $r = .18$), with debate demonstrating the largest effect among communication activities studied. At the individual study level, \textcite{colbert1987effects} found significant critical thinking gains in debaters relative to non-debaters over a full academic year..

Beyond critical thinking, debate participation correlates with substantial academic outcomes. \textcite{mezuk2009urban} found Chicago Debate League participants were 70\% more likely to graduate and three times less likely to drop out than comparable non-participants, with gains concentrated among the lowest-performing students at baseline.

Forensics scholars have proposed theoretical frameworks for these effects, each grounded in established learning theory. \textcite{colbert1995enhancing} argues that critical thinking develops through what he calls ``intellectual disequilibrium,'' a Piagetian concept where confrontation with opposing positions forces learners to restructure their mental models. Debate's adversarial structure ensures this confrontation is unavoidable: a debater cannot ignore the opposing case but must engage with it directly. \textcite{cirlin1997sociological} grounds debate's communicative and socialising effects in Richard Weaver's rhetorical theory and the concept of ``sermonic'' language, while \textcite{hicks1998public}'s ``democratic ethos'' draws on Rawls's ideal of public reason. Drawing these threads together, \textcite{bellon2000research} synthesized evidence from competitive forensics, communication pedagogy, and educational psychology to argue that debate is uniquely effective precisely because it engages cognitive, communicative, social, and civic levels simultaneously. However, these forensics-specific frameworks have not been connected to broader theories of experiential and transformative learning.

% In the Theoretical Framework section of this paper, I apply \textcite{vygotsky1978mind}'s zone of proximal development to the scaffolded learning that occurs through debate coaching and graduated competition, map \textcite{kolb1984experiential}'s experiential learning cycle onto the iterative process of debating, receiving feedback, and refining strategy, and draw on research into structured debriefing in simulation and role-playing \parencite{crookall2014engaging, daniau2016transformative} to illuminate how debate's built-in oral critique enables the reflective processing essential to lasting transformation.

Although forensics pedagogy has documented debate's transformative effects, no prior work analyzes major debate formats from a game design perspective or applies transformative role-playing game theory to competitive debate. Addressing this gap is the central contribution of this thesis.

\subsection{Debate as Role-Playing Game}

In the World Schools format, two teams receive a motion (e.g., ``This House Would legalize the possession of drugs'') and are randomly assigned to Proposition/Government or Opposition. They deliver alternating speeches as if individuals supporting or opposing the motion. Judges listen as someone with no opinion on the motion, expert in debate rules, and averagely informed, deciding who most convinced the character they role-play. The debate ends with the judge ruling who won, explaining the decision, and offering feedback for improvement.

Understanding debate's relationship to role-playing games requires first establishing what role-playing games are and where debate fits among related practices. \textcite{zagal2024definitions} note that defining role-playing games is contentious, but certain characteristics recur: rule-structured creation and enactment of characters in a fictional world, with broad choice over what actions to attempt. \textcite{bowman2025transformative} offer a functional definition: analog role-playing games are ``co-creative experiences in which participants immerse into fictional characters and realities for a bounded period of time and improvise through spontaneous, emergent playfulness'' (p. 16). Debate fits both definitions: it places fictional defenders of positions in a fictional world where a crucial decision will be taken based on their discourse, co-created and time-bound, within which everything is playfully improvised.

Transformative game design literature explicitly positions debate among ``cousin forms'' of analog role-playing games \parencite{bowman2014educational,bowman2025transformative}, alongside simulations, educational role-playing, improv theatre, and psychodrama. The distinction rests on community practices, norms, and innovations, not on the act of role-playing itself. While debate is indeed used in educational contexts, it differs fundamentally from close cousins such as work simulations or classroom exercises in ways that bring it closer to analog role-playing games proper. Three theoretical dimensions illuminate this proximity: the depth of character immersion, the strength of alibi, and the leisure-community dimension.

First, regarding immersion: As opposed to work simulations or boardgames, debate requires players to genuinely inhabit an epistemic stance, constructing and defending arguments as if they believed them. This process of active perspective-taking requires what \textcite{bowman2010functions} describes as adopting a theory of mind, thinking ``as if'' one were someone else in a unique set of circumstances (p.~57). The debater must anticipate objections, identify weaknesses, and respond to challenges as someone who holds the assigned position, a cognitive demand that approaches the immersion into character central to role-playing games. As \textcite{seo2022good} observes, ``the moment one received a motion, one adopted the mindset of a person who was completely convinced of that point of view. One clung to this feeling of absolute conviction to make arguments, sink objections, and display passion'' (p.~222). I and other debaters have observed that the cognitive load of debate performance is incompatible with anything short of fully assuming the assigned opinion for the round's duration. Seo corroborates this from the inside: ``Midspeech, a speaker rarely experiences the physical strain of thinking so damn hard. The adrenaline induced a mental fogginess that numbed the hard edges of experience'' (p.~159). This passionate delivery, especially under such cognitive load, is difficult to sustain through meta-cognitive simulation alone. Affective role-assumption becomes the efficient path, organically eliciting rather than faking the emotions expressed.

\textcite{bowman2010functions} identifies nine types of relationships between players and characters, including the ``Fragmented Self,'' where a facet of the player's personality or interests becomes magnified into a central character trait. Debate characters fit this framework: they represent fragmented selves whose defining feature is epistemic commitment to a particular position. The debater role-plays as someone who believes the motion should pass or fail, adopting that argumentative perspective as the character's core identity. They find in themselves beliefs they can hold that support the position and banish from thought all other, focusing on a fragment of their identity.

Second, regarding alibi: this concept describes how game participants can act with lessened social consequences, encouraging risk-taking and willingness to fail \parencite{deterding2017alibis}. Debate provides robust alibi through its competitive frame. Because debate is ``just a game,'' debaters can advance positions they might not advocate publicly and explore uncomfortable arguments. \textcite{seo2022good} captures this vividly: the random assignment of sides provided a paradoxical ``sense of freedom,'' allowing him to ``flirt with ideas, unencumbered by expectations of consistency or deep conviction'' (p.~20). Work simulations do not offer that, as the situational enactment they demand does not require holding specific beliefs, which are a strong coordinate of identity. This enables the identity experimentation that \textcite{bowman2025transformative} identify as central to transformational containers. Debate, as a competitive game, does not offer alibi to fail intellectually; but many TTRPGs and Nordic larps similarly demand performative or emotional skill while still offering identity alibi. In mainstream North American boffer LARPs, for instance, players who ``break game'' risk derision or ostracism \parencite{brown2016culture}, while in Nordic larp communities, immersion into character is often treated as an indicator of success, with its absence experienced as personal failure \parencite{harder2018larp,pohjola2015steering}. Alibi to fail is not essential to a role-playing game; alibi to assume the mantle of a character is. On that metric, debate offers players a very strong alibi to embody a character with deeply different opinions from their own: what Seo describes as tracing ``the steps of how a sensible person (us!) could arrive at conclusions that might otherwise have seemed alien'' (p.~224).

Third, regarding the leisure-community dimension: competitive debate possesses passionate discourse communities, distinguishing role-playing games from cousin forms. I am part of these communities. The World Schools and British Parliamentary circuits encompass hundreds of competitions. Debaters participate voluntarily, often for decades; they develop specialized terminology, shared references, and community norms; they attend tournaments analogous to gaming conventions. The activity is often pursued for its intrinsic satisfactions, not merely as preparation for external goal. This leisure dimension, combined with the vibrant discourse community, positions competitive debate closer to role-playing games than to the educational simulations with which it shares surface similarities.

Furthermore, evidence for debate's role-playing dimension appears throughout its practice. Experienced coaches regularly employ explicit role-playing techniques during preparation, asking debaters to ``become'' stakeholders affected by policies. Furthermore, debate ethos as I witnessed it is summoned through traditions that explicitly stage parliamentary fiction: bench-pounding, heckling, Points of Information delivered as if holding onto wigs. The format performs a fiction of parliamentary deliberation, with debaters and judges alike participating in theatrical conventions inherited from eighteenth-century British debate culture.

Genuine differences between debate and games that are more traditionally seen as analogue role-playing games, such as Dungeons and Dragons, are undeniable: debate usually lacks the narrative worldbuilding, the character backstories, and the extended campaign structures of tabletop or live-action role-playing. Debate's role-playing operates through the same psychological mechanisms, justifying application of the substantial body of transformative game design scholarship to how design choices might enhance or undermine its educational effects.

\subsection{Literature Review of Format Critiques}

Game design theory addresses problems that forensics pedagogy has not solved. The largest body of work that analyzes debate as an educational tool is that of forensics pedagogy. This centers around the more predominant formats played in the US (especially American Policy). I will thus focus on reviewing what literature was written that echoes the problems I set out to solve and that forms arguments that are also applicable to the formats that I seek to improve (especially the World Schools format).

\subsubsection{The Technical-Public Sphere Tension}

A fundamental critique pervading the literature concerns the evolution of debate formats away from public sphere accessibility toward technical sphere insularity. \textcite{combs2002psychosocial} documents how NDT and CEDA have evolved into technical sphere activities characterized by expert-level topic knowledge, speaking rates exceeding what listeners can process, sophisticated analytical models, and specially trained judges: Colbert found average CEDA debaters speaking at approximately 237 words per minute and NDT finalists at 284, rates that ``greatly exceed what is considered optimal for normal speaking'' \parencite[p.~36]{combs2002psychosocial}.

While World Schools and British Parliamentary formats explicitly discourage fast delivery through style-based evaluation, similar pressures toward information density persist. These formats evaluate on two co-equal axes, Style and Content, yet increasing information density improves argumentative quality while reducing intelligibility. Because argumentative quality is perceived as more objective than style in an academic setting, the competitive meta gradually shifts toward privileging Content over Style despite the rubric formally treating them as equal. The result is that skills developed may not transfer to contexts without expert judges, a difficulty captured in the audience-standard literature, where \textcite{hicks2002public} and \textcite{weiss2002audience} disagree on whether debaters should be persuasive to a regular audience or to the perfect rational listener, with educational value lost either way.

\subsubsection{Value Convergence and Ideological Homogenization}

The ``closed system'' dynamic described by Harris and Smith, cited in \textcite{combs2002psychosocial}, creates conditions conducive to value convergence. In academic debate, tournament practices are insulated from external feedback and where judges who deviate from community norms are ``effectively excluded from the system much like the human system rejects certain bacteria'' \parencite[p.~40]{combs2002psychosocial}. This dynamic may cultivate shared assumptions that, while facilitating intra-community communication, potentially diminish debaters' capacity to engage persuasively with those holding different values.

Crenshaw's work on feminism in debate, cited by \textcite{mcgee2002emancipatory}, illustrates how ``argument borrowing'' can lead to reductionist treatment of complex intellectual traditions. Debaters observe which arguments succeed in out-rounds, then adopt those lines of reasoning for their strategic value rather than out of genuine engagement with the underlying ideas, producing superficial understanding and reduced capacity for authentic cross-ideological dialogue.


\subsubsection{Hierarchical Dynamics and Psychological Safety}

Forensics pedagogy is thinnest on this issue, yet several authors recognize the risk. \textcite{snider1999code} addresses safety and ethics in his Code of the Debate, and in his later practical guide warns organizers to ``avoid a counterproductive rigid and hierarchical organization,'' noting ``I've seen groups die when they let ego and status get in the way'' \parencite[p.~36]{snider2014sparking}. The International Debate Education Association's ethics guide for coaches and youth workers similarly acknowledges the ``unbalanced power dynamic'' inherent in coach-student relationships, compounded by the fact that ``older debaters often mentor younger teams'' and thus occupy dual roles of peer and authority figure \parencite{idea2024ethics}. Riley and Hollihan further suggest that approximately three years of debate experience are required to compete effectively on the national circuit, effectively excluding students without extensive preparation and reinforcing experience-based hierarchies \parencite[p.~38]{combs2002psychosocial}.

The mechanism behind these dynamics is illuminated by recent research. \textcite{prussikora2025cracovian}, citing Simon et al.\ (2020), note that ``people placed in competitive intellectual situations tend to develop an adversarial mindset'' characterized by myside bias and otherside bias, distorting cognitive appraisal of the issues being debated. They identify competitive debate's ``most crucial weakness'': its rules ``do not allow for saying to the opponent: `You are right, I was wrong'\,'' instead training debaters to listen for omissions and inaccuracies rather than for understanding \parencite[p.~1474]{prussikora2025cracovian}. When win-loss records become the primary measure of status within these hermetic communities, the adversarial structure produces natural hierarchies in which winning debaters gain social authority while losing debaters are marginalized, reinforcing competitive pressure at the expense of collaborative learning.

\subsection{Gap Identification}

This literature review reveals a significant gap. Forensics pedagogy has established debate's educational benefits but has not applied game design frameworks to analyze how format mechanics produce specific effects; transformative role-playing game theory has not been applied to competitive debate. Addressing this gap is a central contribution of this thesis.
